\documentclass[multi=circuitikz,border=10pt]{standalone}
\usepackage[utf8]{inputenc}
\usepackage[T1]{fontenc}
\usepackage{amsmath}
\usepackage{circuitikz}
\ctikzset{inductors/scale=0.8}
\tikzset{bus/.style={line width=2.2pt}}

\newcommand{\refline}{\draw[thick] (-1,0) -- (11.8,0) node[right]{Ref.};}
\newcommand{\bus}[2]{\draw[bus] (#1,2.1) -- (#1,2.9); \node[font=\scriptsize] at (#1,3.3) {#2};}
\newcommand{\linhaZero}{\draw (6.6,2.5) to[L, l={$Z_{0,L} = j0{,}2666$ pu}] (9.0,2.5) -- (9.6,2.5);
  \bus{9.6}{F}\bus{6.6}{T (245 kV)}
  \draw[dashed, gray] (9.6,2.1) -- (9.6,0); \draw (9.6,2.1) to[open, v^=$V_b^{(0)}$] (9.6,0);}
\newcommand{\geradorZero}{\draw (0,0) to[L, l_={$3Z_n$}] (0,1.3) to[L, l_={$X_{g0} = j0{,}0125$ pu}] (0,2.5) -- (3.6,2.5); \bus{3.6}{G (15 kV)}}
\newcommand{\aberto}[2]{\draw (#1,2.5) -- (#1+0.7,2.5); \fill (#1+0.7,2.5) circle (1.6pt);
  \draw (#1+0.7,2.5) -- (#1+1.7,2.95); \fill (#1+2.0,2.5) circle (1.6pt); \draw (#1+2.0,2.5) -- (#2,2.5);}

\begin{document}
% ---------- 1) Delta-Delta ----------
\begin{circuitikz}[american, font=\footnotesize]
  \refline
  \draw (0,0) -- (0,2.5) -- (0.6,2.5) to[L, l={$X_{g0} = j0{,}0125$ pu}] (3.0,2.5) -- (3.6,2.5); \bus{3.6}{G (15 kV)}
  \aberto{3.6}{6.6}
  \node[font=\scriptsize, align=center] at (5.1,3.9) {$\Delta$-$\Delta$: circuito aberto};
  \node[font=\scriptsize] at (5.1,1.8) {($X_t$ não conduz $I_0$)};
  \linhaZero
  \node[font=\small] at (5,-1.1) {$\Delta$-$\Delta$: $Z_0$ visto de F $= \infty$ $\Rightarrow$ $I_{1\phi} = 0$};
\end{circuitikz}

% ---------- 2) Yg-Yg ----------
\begin{circuitikz}[american, font=\footnotesize]
  \refline
  \geradorZero
  \draw (3.6,2.5) to[L, l={$X_t = j0{,}03$ pu}] (6.0,2.5) -- (6.6,2.5);
  \linhaZero
  \node[font=\small, align=center] at (5,-1.2) {Yg–Yg: $Z_0 = X_{g0} + 3Z_n + X_t + Z_{0,L}$ (ex.: $Z_n = 0$ $\Rightarrow$ $j0{,}3091$ pu)\\
   $I_{1\phi} = 3/(Z_0+Z_1+Z_2) = 5{,}06$ pu $\approx 1194$ A};
\end{circuitikz}

% ---------- 3) Delta (gerador) - Yg (linha) ----------
\begin{circuitikz}[american, font=\footnotesize]
  \refline
  \bus{3.6}{G (15 kV)}
  \draw[gray, dashed] (0.5,2.5) -- (3.6,2.5);
  \node[font=\scriptsize, gray, align=center] at (1.4,4.0) {gerador ($X_{g0}$, $3Z_n$)\\isolado pelo $\Delta$};
  \draw (6.6,2.5) to[L, l={$X_t = j0{,}03$ pu}] (6.6,0);
  \linhaZero
  \node[font=\small, align=center] at (5,-1.2) {$\Delta$–Yg: $Z_0 = X_t + Z_{0,L} = j0{,}2966$ pu (independe do gerador e de $Z_n$)\\
   $I_{1\phi} = 5{,}17$ pu $\approx 1219$ A};
\end{circuitikz}

% ---------- 4) Yg (gerador) - Delta (linha) ----------
\begin{circuitikz}[american, font=\footnotesize]
  \refline
  \geradorZero
  \draw (3.6,2.5) to[L, l={$X_t = j0{,}03$ pu}] (3.6,0);
  \linhaZero
  \node[font=\scriptsize, align=center] at (5.1,3.9) {lado da linha em $\Delta$:\\barra T aberta};
  \node[font=\small] at (5,-1.1) {Yg–$\Delta$: $Z_0$ visto de F $= \infty$ $\Rightarrow$ $I_{1\phi} = 0$};
\end{circuitikz}

% ---------- 5) Y isolado (gerador) - Yg (linha) ----------
\begin{circuitikz}[american, font=\footnotesize]
  \refline
  \draw (0,0) -- (0,2.5) -- (0.6,2.5) to[L, l={$X_{g0}$}] (3.0,2.5) -- (3.6,2.5); \bus{3.6}{G (15 kV)}
  \aberto{3.6}{6.6}
  \node[font=\scriptsize, align=center] at (5.1,3.9) {neutro do gerador-lado Y isolado:\\sem corrente de seq. zero (caso 2)};
  \linhaZero
  \node[font=\small] at (5,-1.1) {Y isolado–Yg: $Z_0 = \infty$ $\Rightarrow$ $I_{1\phi} = 0$};
\end{circuitikz}
\end{document}
